% Options for packages loaded elsewhere
\PassOptionsToPackage{unicode}{hyperref}
\PassOptionsToPackage{hyphens}{url}
\PassOptionsToPackage{dvipsnames,svgnames,x11names}{xcolor}
%
\documentclass[
  12pt,
]{book}
\usepackage{amsmath,amssymb}
\usepackage{setspace}
\usepackage{iftex}
\ifPDFTeX
  \usepackage[T1]{fontenc}
  \usepackage[utf8]{inputenc}
  \usepackage{textcomp} % provide euro and other symbols
\else % if luatex or xetex
  \usepackage{unicode-math} % this also loads fontspec
  \defaultfontfeatures{Scale=MatchLowercase}
  \defaultfontfeatures[\rmfamily]{Ligatures=TeX,Scale=1}
\fi
\usepackage{lmodern}
\ifPDFTeX\else
  % xetex/luatex font selection
\fi
% Use upquote if available, for straight quotes in verbatim environments
\IfFileExists{upquote.sty}{\usepackage{upquote}}{}
\IfFileExists{microtype.sty}{% use microtype if available
  \usepackage[]{microtype}
  \UseMicrotypeSet[protrusion]{basicmath} % disable protrusion for tt fonts
}{}
\makeatletter
\@ifundefined{KOMAClassName}{% if non-KOMA class
  \IfFileExists{parskip.sty}{%
    \usepackage{parskip}
  }{% else
    \setlength{\parindent}{0pt}
    \setlength{\parskip}{6pt plus 2pt minus 1pt}}
}{% if KOMA class
  \KOMAoptions{parskip=half}}
\makeatother
\usepackage{xcolor}
\usepackage[left=4cm, right=3cm, top=2.5cm, bottom=2.5cm]{geometry}
\usepackage{longtable,booktabs,array}
\usepackage{calc} % for calculating minipage widths
% Correct order of tables after \paragraph or \subparagraph
\usepackage{etoolbox}
\makeatletter
\patchcmd\longtable{\par}{\if@noskipsec\mbox{}\fi\par}{}{}
\makeatother
% Allow footnotes in longtable head/foot
\IfFileExists{footnotehyper.sty}{\usepackage{footnotehyper}}{\usepackage{footnote}}
\makesavenoteenv{longtable}
\usepackage{graphicx}
\makeatletter
\def\maxwidth{\ifdim\Gin@nat@width>\linewidth\linewidth\else\Gin@nat@width\fi}
\def\maxheight{\ifdim\Gin@nat@height>\textheight\textheight\else\Gin@nat@height\fi}
\makeatother
% Scale images if necessary, so that they will not overflow the page
% margins by default, and it is still possible to overwrite the defaults
% using explicit options in \includegraphics[width, height, ...]{}
\setkeys{Gin}{width=\maxwidth,height=\maxheight,keepaspectratio}
% Set default figure placement to htbp
\makeatletter
\def\fps@figure{htbp}
\makeatother
\setlength{\emergencystretch}{3em} % prevent overfull lines
\providecommand{\tightlist}{%
  \setlength{\itemsep}{0pt}\setlength{\parskip}{0pt}}
\setcounter{secnumdepth}{5}
\usepackage[utf8]{inputenc}
\usepackage[spanish,es-tabla]{babel}
\usepackage[fixlanguage]{babelbib}
\usepackage{geometry}
\geometry{letterpaper,left=2cm,top=2cm, right=2cm}
\usepackage{times}           
\usepackage{caption}
\captionsetup[figure, table]{labelfont={bf},labelformat={default},labelsep=period}
\usepackage{graphicx}
\usepackage{booktabs}
\usepackage{longtable}
\usepackage{array}
\usepackage{multirow}
\usepackage{wrapfig}
%\usepackage{float}
\usepackage{colortbl}
\usepackage{xcolor}
\usepackage{pdflscape}
\usepackage{tabu}
\usepackage{threeparttable}
\usepackage{pdfpages}
\usepackage{hyphsubst}
\usepackage{floatrow}
\floatsetup[figure]{capposition=top}
\floatsetup[table]{capposition=top}
\usepackage{booktabs}
\usepackage{longtable}
\usepackage{array}
\usepackage{multirow}
\usepackage{wrapfig}
\usepackage{float}
\usepackage{colortbl}
\usepackage{pdflscape}
\usepackage{tabu}
\usepackage{threeparttable}
\usepackage{threeparttablex}
\usepackage[normalem]{ulem}
\usepackage{makecell}
\usepackage{xcolor}
\ifLuaTeX
  \usepackage{selnolig}  % disable illegal ligatures
\fi
\IfFileExists{bookmark.sty}{\usepackage{bookmark}}{\usepackage{hyperref}}
\IfFileExists{xurl.sty}{\usepackage{xurl}}{} % add URL line breaks if available
\urlstyle{same}
\hypersetup{
  pdftitle={Efectos políticos de las huelgas},
  pdfauthor={Equipo Fondecyt},
  colorlinks=true,
  linkcolor={blue},
  filecolor={Maroon},
  citecolor={Blue},
  urlcolor={Blue},
  pdfcreator={LaTeX via pandoc}}

\title{Efectos políticos de las huelgas}
\usepackage{etoolbox}
\makeatletter
\providecommand{\subtitle}[1]{% add subtitle to \maketitle
  \apptocmd{\@title}{\par {\large #1 \par}}{}{}
}
\makeatother
\subtitle{con ELSOC 2016-2022}
\author{Equipo Fondecyt}
\date{}

\begin{document}
\maketitle

{
\hypersetup{linkcolor=}
\setcounter{tocdepth}{1}
\tableofcontents
}
\listoffigures
\listoftables
\setstretch{1.5}
\hypertarget{presentaciuxf3n}{%
\chapter*{Presentación}\label{presentaciuxf3n}}
\addcontentsline{toc}{chapter}{Presentación}

Preparación y exploración de datos ELSOC 2016-2022 para analizar variables que puedan relacionarse a los efectos políticos de las huelgas.

Se debe considerar que los ítems del cuestionario ELSOC no se presentan en todas las olas.

Las olas corresponden a los siguientes años:

\begin{itemize}
\tightlist
\item
  Ola 1: 2016
\item
  Ola 2: 2017
\item
  Ola 3: 2018
\item
  Ola 4: 2019
\item
  Ola 5: 2021
\item
  Ola 6: 2022
\end{itemize}

En mayo de 2024 se contará con la séptima ola correspondiente al año 2023.

\hypertarget{preparaciuxf3n-y-exploraciuxf3n-datos}{%
\chapter{Preparación y exploración datos}\label{preparaciuxf3n-y-exploraciuxf3n-datos}}

\hypertarget{exploraciuxf3n-variables-de-interuxe9s}{%
\section{Exploración variables de interés}\label{exploraciuxf3n-variables-de-interuxe9s}}

Primero, una revisión exploratoria de las variables de interés, considerando todas las olas de la base de datos y la muestra original y de refresco, sin ponderar.

\hypertarget{caracterizaciuxf3n-empleo}{%
\subsection{Caracterización empleo}\label{caracterizaciuxf3n-empleo}}

\hypertarget{actividad}{%
\subsubsection{Actividad}\label{actividad}}

\begin{itemize}
\tightlist
\item
  m02. ¿Cuál de estas situaciones describe mejor su actividad principal durante el último mes?
\end{itemize}

Frecuencias y porcentajes de Actividad, por ola

Actividad

Ola

Total

2016

2017

2018

2019

2021

2022

Trabaja de maneraremunerada conjornada completa

{1391}{47.6~\%}

{1116}{45.2~\%}

{1688}{45.2~\%}

{1501}{44.1~\%}

{1095}{40~\%}

{1137}{41.8~\%}

{7928}{44.1~\%}

Trabaja de maneraremunerada a tiempoparcial

{351}{12~\%}

{362}{14.7~\%}

{537}{14.4~\%}

{560}{16.4~\%}

{484}{17.7~\%}

{456}{16.7~\%}

{2750}{15.3~\%}

Estudia y trabaja

{60}{2.1~\%}

{71}{2.9~\%}

{90}{2.4~\%}

{68}{2~\%}

{62}{2.3~\%}

{49}{1.8~\%}

{400}{2.2~\%}

Solo estudia

{101}{3.5~\%}

{75}{3~\%}

{116}{3.1~\%}

{77}{2.3~\%}

{51}{1.9~\%}

{36}{1.3~\%}

{456}{2.5~\%}

Jubilado opensionado

{371}{12.7~\%}

{345}{14~\%}

{479}{12.8~\%}

{499}{14.6~\%}

{433}{15.8~\%}

{545}{20~\%}

{2672}{14.9~\%}

Desempleado,buscando trabajo

{166}{5.7~\%}

{73}{3~\%}

{198}{5.3~\%}

{154}{4.5~\%}

{172}{6.3~\%}

{119}{4.4~\%}

{882}{4.9~\%}

Realiza tareas noremuneradas

{367}{12.6~\%}

{332}{13.5~\%}

{443}{11.9~\%}

{449}{13.2~\%}

{366}{13.4~\%}

{294}{10.8~\%}

{2251}{12.5~\%}

Esta enfermo o tieneuna discapacidad

{24}{0.8~\%}

{17}{0.7~\%}

{29}{0.8~\%}

{32}{0.9~\%}

{44}{1.6~\%}

{28}{1~\%}

{174}{1~\%}

No estudia, notrabaja y no buscatrabajo

{90}{3.1~\%}

{77}{3.1~\%}

{152}{4.1~\%}

{67}{2~\%}

{33}{1.2~\%}

{59}{2.2~\%}

{478}{2.7~\%}

Total

{2921}{100~\%}

{2468}{100~\%}

{3732}{100~\%}

{3407}{100~\%}

{2740}{100~\%}

{2723}{100~\%}

{17991}{100~\%}

\hypertarget{ocupaciuxf3n}{%
\subsubsection{Ocupación}\label{ocupaciuxf3n}}

\begin{itemize}
\tightlist
\item
  m03. ¿Cuál es su ocupación u oficio actual? Describa sus principales tareas y funciones en el puesto de trabajo actual
\end{itemize}

Esta variable se encuentra codificada de acuerdo con la CIUO88 para el año 2016, y CIUO08 para los años 2018 y 2021.

\begin{itemize}
\tightlist
\item
  ciuo88\_m03: CIUO (1988) del entrevistado
\item
  ciuo08\_m03: CIUO (2008) del entrevistado
\end{itemize}

Por lo tanto, se realiza una recodificación generando una nueva variable donde están todas las olas homolagadas a la CIUO08. Se presentan las frecuecnias y porcentajes con uno y dos dígitos.

Frecuencias y porcentajes de Ocupación con 1 dígito, por ola

CIUO08

Ola

Total

2016

2017

2018

2019

2021

2022

1

{89}{5~\%}

{0}{0~\%}

{58}{2.5~\%}

{0}{0~\%}

{41}{2.5~\%}

{0}{0~\%}

{188}{3.3~\%}

2

{257}{14.4~\%}

{0}{0~\%}

{308}{13.3~\%}

{0}{0~\%}

{248}{15.1~\%}

{0}{0~\%}

{813}{14.2~\%}

3

{362}{20.3~\%}

{0}{0~\%}

{175}{7.6~\%}

{0}{0~\%}

{147}{9~\%}

{0}{0~\%}

{684}{11.9~\%}

4

{26}{1.5~\%}

{0}{0~\%}

{187}{8.1~\%}

{0}{0~\%}

{127}{7.7~\%}

{0}{0~\%}

{340}{5.9~\%}

5

{348}{19.5~\%}

{0}{0~\%}

{573}{24.8~\%}

{0}{0~\%}

{378}{23~\%}

{0}{0~\%}

{1299}{22.6~\%}

6

{41}{2.3~\%}

{0}{0~\%}

{25}{1.1~\%}

{0}{0~\%}

{30}{1.8~\%}

{0}{0~\%}

{96}{1.7~\%}

7

{319}{17.9~\%}

{0}{0~\%}

{343}{14.8~\%}

{0}{0~\%}

{294}{17.9~\%}

{0}{0~\%}

{956}{16.7~\%}

8

{120}{6.7~\%}

{0}{0~\%}

{202}{8.7~\%}

{0}{0~\%}

{125}{7.6~\%}

{0}{0~\%}

{447}{7.8~\%}

9

{220}{12.3~\%}

{0}{0~\%}

{444}{19.2~\%}

{0}{0~\%}

{251}{15.3~\%}

{0}{0~\%}

{915}{15.9~\%}

Total

{1782}{100~\%}

{0}{100~\%}

{2315}{100~\%}

{0}{100~\%}

{1641}{100~\%}

{0}{100~\%}

{5738}{100~\%}

Frecuencias y porcentajes de Ocupación con 2 dígitos, por ola

CIUO08

Ola

Total

2016

2017

2018

2019

2021

2022

11

{0}{0~\%}

{0}{0~\%}

{6}{0.3~\%}

{0}{0~\%}

{8}{0.5~\%}

{0}{0~\%}

{14}{0.2~\%}

12

{14}{0.8~\%}

{0}{0~\%}

{8}{0.3~\%}

{0}{0~\%}

{2}{0.1~\%}

{0}{0~\%}

{24}{0.4~\%}

13

{15}{0.8~\%}

{0}{0~\%}

{22}{1~\%}

{0}{0~\%}

{3}{0.2~\%}

{0}{0~\%}

{40}{0.7~\%}

14

{60}{3.4~\%}

{0}{0~\%}

{22}{1~\%}

{0}{0~\%}

{28}{1.7~\%}

{0}{0~\%}

{110}{1.9~\%}

21

{47}{2.6~\%}

{0}{0~\%}

{52}{2.2~\%}

{0}{0~\%}

{50}{3~\%}

{0}{0~\%}

{149}{2.6~\%}

22

{34}{1.9~\%}

{0}{0~\%}

{37}{1.6~\%}

{0}{0~\%}

{26}{1.6~\%}

{0}{0~\%}

{97}{1.7~\%}

23

{82}{4.6~\%}

{0}{0~\%}

{125}{5.4~\%}

{0}{0~\%}

{86}{5.2~\%}

{0}{0~\%}

{293}{5.1~\%}

24

{45}{2.5~\%}

{0}{0~\%}

{41}{1.8~\%}

{0}{0~\%}

{35}{2.1~\%}

{0}{0~\%}

{121}{2.1~\%}

25

{6}{0.3~\%}

{0}{0~\%}

{7}{0.3~\%}

{0}{0~\%}

{15}{0.9~\%}

{0}{0~\%}

{28}{0.5~\%}

26

{43}{2.4~\%}

{0}{0~\%}

{46}{2~\%}

{0}{0~\%}

{36}{2.2~\%}

{0}{0~\%}

{125}{2.2~\%}

31

{54}{3~\%}

{0}{0~\%}

{32}{1.4~\%}

{0}{0~\%}

{37}{2.3~\%}

{0}{0~\%}

{123}{2.1~\%}

32

{29}{1.6~\%}

{0}{0~\%}

{42}{1.8~\%}

{0}{0~\%}

{42}{2.6~\%}

{0}{0~\%}

{113}{2~\%}

33

{209}{11.7~\%}

{0}{0~\%}

{68}{2.9~\%}

{0}{0~\%}

{54}{3.3~\%}

{0}{0~\%}

{331}{5.8~\%}

34

{63}{3.5~\%}

{0}{0~\%}

{24}{1~\%}

{0}{0~\%}

{9}{0.5~\%}

{0}{0~\%}

{96}{1.7~\%}

35

{7}{0.4~\%}

{0}{0~\%}

{8}{0.3~\%}

{0}{0~\%}

{5}{0.3~\%}

{0}{0~\%}

{20}{0.3~\%}

36

{0}{0~\%}

{0}{0~\%}

{1}{0~\%}

{0}{0~\%}

{0}{0~\%}

{0}{0~\%}

{1}{0~\%}

41

{2}{0.1~\%}

{0}{0~\%}

{74}{3.2~\%}

{0}{0~\%}

{11}{0.7~\%}

{0}{0~\%}

{87}{1.5~\%}

42

{24}{1.3~\%}

{0}{0~\%}

{43}{1.9~\%}

{0}{0~\%}

{24}{1.5~\%}

{0}{0~\%}

{91}{1.6~\%}

43

{0}{0~\%}

{0}{0~\%}

{57}{2.5~\%}

{0}{0~\%}

{55}{3.4~\%}

{0}{0~\%}

{112}{2~\%}

44

{0}{0~\%}

{0}{0~\%}

{13}{0.6~\%}

{0}{0~\%}

{36}{2.2~\%}

{0}{0~\%}

{49}{0.9~\%}

47

{0}{0~\%}

{0}{0~\%}

{0}{0~\%}

{0}{0~\%}

{1}{0.1~\%}

{0}{0~\%}

{1}{0~\%}

51

{54}{3~\%}

{0}{0~\%}

{134}{5.8~\%}

{0}{0~\%}

{84}{5.1~\%}

{0}{0~\%}

{272}{4.7~\%}

52

{180}{10.1~\%}

{0}{0~\%}

{299}{12.9~\%}

{0}{0~\%}

{182}{11.1~\%}

{0}{0~\%}

{661}{11.5~\%}

53

{49}{2.7~\%}

{0}{0~\%}

{76}{3.3~\%}

{0}{0~\%}

{65}{4~\%}

{0}{0~\%}

{190}{3.3~\%}

54

{65}{3.6~\%}

{0}{0~\%}

{64}{2.8~\%}

{0}{0~\%}

{47}{2.9~\%}

{0}{0~\%}

{176}{3.1~\%}

61

{34}{1.9~\%}

{0}{0~\%}

{22}{1~\%}

{0}{0~\%}

{19}{1.2~\%}

{0}{0~\%}

{75}{1.3~\%}

62

{7}{0.4~\%}

{0}{0~\%}

{1}{0~\%}

{0}{0~\%}

{10}{0.6~\%}

{0}{0~\%}

{18}{0.3~\%}

63

{0}{0~\%}

{0}{0~\%}

{2}{0.1~\%}

{0}{0~\%}

{1}{0.1~\%}

{0}{0~\%}

{3}{0.1~\%}

71

{123}{6.9~\%}

{0}{0~\%}

{109}{4.7~\%}

{0}{0~\%}

{96}{5.9~\%}

{0}{0~\%}

{328}{5.7~\%}

72

{62}{3.5~\%}

{0}{0~\%}

{82}{3.5~\%}

{0}{0~\%}

{51}{3.1~\%}

{0}{0~\%}

{195}{3.4~\%}

73

{7}{0.4~\%}

{0}{0~\%}

{20}{0.9~\%}

{0}{0~\%}

{11}{0.7~\%}

{0}{0~\%}

{38}{0.7~\%}

74

{29}{1.6~\%}

{0}{0~\%}

{34}{1.5~\%}

{0}{0~\%}

{31}{1.9~\%}

{0}{0~\%}

{94}{1.6~\%}

75

{98}{5.5~\%}

{0}{0~\%}

{98}{4.2~\%}

{0}{0~\%}

{105}{6.4~\%}

{0}{0~\%}

{301}{5.2~\%}

81

{6}{0.3~\%}

{0}{0~\%}

{52}{2.2~\%}

{0}{0~\%}

{37}{2.3~\%}

{0}{0~\%}

{95}{1.7~\%}

82

{0}{0~\%}

{0}{0~\%}

{1}{0~\%}

{0}{0~\%}

{1}{0.1~\%}

{0}{0~\%}

{2}{0~\%}

83

{114}{6.4~\%}

{0}{0~\%}

{149}{6.4~\%}

{0}{0~\%}

{87}{5.3~\%}

{0}{0~\%}

{350}{6.1~\%}

91

{154}{8.6~\%}

{0}{0~\%}

{206}{8.9~\%}

{0}{0~\%}

{152}{9.3~\%}

{0}{0~\%}

{512}{8.9~\%}

92

{15}{0.8~\%}

{0}{0~\%}

{43}{1.9~\%}

{0}{0~\%}

{36}{2.2~\%}

{0}{0~\%}

{94}{1.6~\%}

93

{27}{1.5~\%}

{0}{0~\%}

{83}{3.6~\%}

{0}{0~\%}

{36}{2.2~\%}

{0}{0~\%}

{146}{2.5~\%}

94

{0}{0~\%}

{0}{0~\%}

{24}{1~\%}

{0}{0~\%}

{5}{0.3~\%}

{0}{0~\%}

{29}{0.5~\%}

95

{14}{0.8~\%}

{0}{0~\%}

{18}{0.8~\%}

{0}{0~\%}

{15}{0.9~\%}

{0}{0~\%}

{47}{0.8~\%}

96

{10}{0.6~\%}

{0}{0~\%}

{16}{0.7~\%}

{0}{0~\%}

{7}{0.4~\%}

{0}{0~\%}

{33}{0.6~\%}

99

{0}{0~\%}

{0}{0~\%}

{54}{2.3~\%}

{0}{0~\%}

{0}{0~\%}

{0}{0~\%}

{54}{0.9~\%}

Total

{1782}{100~\%}

{0}{100~\%}

{2315}{100~\%}

{0}{100~\%}

{1641}{100~\%}

{0}{100~\%}

{5738}{100~\%}

\hypertarget{nivel-de-cualificaciuxf3n}{%
\subsubsection{Nivel de cualificación}\label{nivel-de-cualificaciuxf3n}}

Se crea una nueva variable con el Nivel de cualificación, de acuerdo con la variable CIUO08 de 2 dígitos \emph{isco08\_2d}, siguiendo el criterio:

\begin{itemize}
\tightlist
\item
  Expertos (Experts): 11 a 34 con universitaria completa o más
\item
  Calificados (Skilled): 11 a 34 con universitaria incompleta o menos, y 35,41,42,43,44,53, y 51,54,61,62,71,72,73,74,75,81 con media completa o más.
\item
  No calificados (Unskilled): 51,54,61,62,71,72,73,74,75,81 con media incompleta o menos, y 52,61,63,75,82,83,90+
\end{itemize}

Frecuencias y porcentajes del Nivel de Cualificación, por ola

Nivel decualificación

Ola

Total

2016

2017

2018

2019

2021

2022

Experts

{240}{15.4~\%}

{0}{0~\%}

{289}{15.5~\%}

{0}{0~\%}

{242}{17.4~\%}

{0}{0~\%}

{771}{16~\%}

Skilled

{650}{41.6~\%}

{0}{0~\%}

{671}{35.9~\%}

{0}{0~\%}

{510}{36.7~\%}

{0}{0~\%}

{1831}{38~\%}

Unskilled

{672}{43~\%}

{0}{0~\%}

{910}{48.7~\%}

{0}{0~\%}

{637}{45.9~\%}

{0}{0~\%}

{2219}{46~\%}

Total

{1562}{100~\%}

{0}{100~\%}

{1870}{100~\%}

{0}{100~\%}

{1389}{100~\%}

{0}{100~\%}

{4821}{100~\%}

\hypertarget{rama-de-actividad-econuxf3mica}{%
\subsubsection{Rama de Actividad Económica}\label{rama-de-actividad-econuxf3mica}}

\begin{itemize}
\tightlist
\item
  m04. ¿A qué rubro o giro de actividad se dedica principalmente la empresa, institución o negocio para la cual usted trabaja?
\end{itemize}

Esta variable se encuentra codificada de acuerdo con la CIIU3 para el año 2016, y CIIU4 para los años 2018 y 2021.

\begin{itemize}
\tightlist
\item
  ciiu3\_m04: CIIU (III version) del entrevistado
\item
  ciiu4\_m04: CIIU (IV version) del entrevistado
\end{itemize}

Por lo tanto, se realiza una recodificación generando una nueva variable donde están todas las olas homolagadas a la CIIU4.

\begin{verbatim}
## 
##   1   2   5  11  13  14  15  17  18  19  20  21  22  24  25  26  27  28  29  32 
##  54  20   7   1  26   2  62   5  31   3  12   6   8   6   2   2   7  14   8   1 
##  35  36  37  40  41  45  50  51  52  55  60  61  63  64  65  66  67  70  71  72 
##   1  18   1   1   3 185  54  16 249  74  80   6  10  23  14  11   2  18   5   8 
##  73  74  75  80  85  90  91  92  93  95  99 
##   2 109  62 161  83   3   6  29  24 150  25
\end{verbatim}

\begin{verbatim}
## 
##   1   2   3   4   5   7   8  10  11  12  13  14  15  16  17  18  19  20  21  22 
##  93  34  22  41   1   4  12 152   7   2  15  54   5  28  17  15   2  13   8  11 
##  23  24  25  26  28  30  31  32  33  34  35  36  38  40  41  42  43  45  46  47 
##   5  13  25   5  10   2  27   5  18   1  18  10   6   7 227  31 100  82  66 560 
##  49  50  51  52  53  55  56  58  59  60  61  62  63  64  65  66  68  69  70  71 
## 180  10   5  45   5  17 179   2   2   5  32  20   1  29  13   9  24  41  12  24 
##  72  73  74  75  77  78  79  80  81  82  84  85  86  87  88  90  91  92  93  94 
##   7  18   6   6   9  12   3  54  87  14 203 358 205  20   6  15   2   4  19  15 
##  95  96  97  99 
##  24  71 300  89
\end{verbatim}

\hypertarget{supervisiuxf3n}{%
\subsubsection{Supervisión}\label{supervisiuxf3n}}

\begin{itemize}
\tightlist
\item
  m06. En su trabajo, ¿a cuántas personas supervisa usted?
\end{itemize}

Se recodifica en una nueva variable donde:

\begin{itemize}
\tightlist
\item
  1 = supervisa a una o más personas
\item
  0 = no supervisa
\end{itemize}

Frecuencias y porcentajes de variable Supervisa, por ola

Supervisa

Ola

Total

2016

2017

2018

2019

2021

2022

0

{1339}{74.3~\%}

{0}{0~\%}

{1769}{77~\%}

{0}{0~\%}

{750}{61.5~\%}

{0}{0~\%}

{3858}{72.5~\%}

1

{463}{25.7~\%}

{0}{0~\%}

{528}{23~\%}

{0}{0~\%}

{469}{38.5~\%}

{0}{0~\%}

{1460}{27.5~\%}

Total

{1802}{100~\%}

{0}{100~\%}

{2297}{100~\%}

{0}{100~\%}

{1219}{100~\%}

{0}{100~\%}

{5318}{100~\%}

\hypertarget{relaciuxf3n-de-empleo}{%
\subsubsection{Relación de empleo}\label{relaciuxf3n-de-empleo}}

\begin{itemize}
\tightlist
\item
  m07. En su actual ocupación usted trabaja como:
\end{itemize}

Frecuencias y porcentajes de la Relación de empleo, por ola

Relación de empleo

Ola

Total

2016

2017

2018

2019

2021

2022

Empleado u obrero enempresa privada

{1091}{61.1~\%}

{0}{0~\%}

{1338}{58.4~\%}

{0}{0~\%}

{883}{53.9~\%}

{0}{0~\%}

{3312}{57.9~\%}

Empleado u obrerodel sector público

{186}{10.4~\%}

{0}{0~\%}

{290}{12.7~\%}

{0}{0~\%}

{228}{13.9~\%}

{0}{0~\%}

{704}{12.3~\%}

Miembro de lasFuerzas Armadas y deOrden

{25}{1.4~\%}

{0}{0~\%}

{25}{1.1~\%}

{0}{0~\%}

{18}{1.1~\%}

{0}{0~\%}

{68}{1.2~\%}

Patrón/a oempleador/a

{85}{4.8~\%}

{0}{0~\%}

{105}{4.6~\%}

{0}{0~\%}

{86}{5.3~\%}

{0}{0~\%}

{276}{4.8~\%}

Trabaja solo, notiene empleados

{321}{18~\%}

{0}{0~\%}

{409}{17.9~\%}

{0}{0~\%}

{345}{21.1~\%}

{0}{0~\%}

{1075}{18.8~\%}

Familiar noremunerado

{4}{0.2~\%}

{0}{0~\%}

{6}{0.3~\%}

{0}{0~\%}

{1}{0.1~\%}

{0}{0~\%}

{11}{0.2~\%}

Servicio doméstico

{75}{4.2~\%}

{0}{0~\%}

{118}{5.2~\%}

{0}{0~\%}

{77}{4.7~\%}

{0}{0~\%}

{270}{4.7~\%}

Total

{1787}{100~\%}

{0}{100~\%}

{2291}{100~\%}

{0}{100~\%}

{1638}{100~\%}

{0}{100~\%}

{5716}{100~\%}

\hypertarget{temporalidad-del-trabajo}{%
\subsubsection{Temporalidad del trabajo}\label{temporalidad-del-trabajo}}

\begin{itemize}
\tightlist
\item
  m08. Su trabajo o negocio principal es de tipo:
\end{itemize}

Frecuencias y porcentajes de la Temporalidad del trabajo, por ola

Temporalidad

Ola

Total

2016

2017

2018

2019

2021

2022

Permanente

{1465}{81.5~\%}

{0}{0~\%}

{1766}{76.6~\%}

{0}{0~\%}

{1203}{73.4~\%}

{0}{0~\%}

{4434}{77.2~\%}

De temporada oestacional

{141}{7.8~\%}

{0}{0~\%}

{228}{9.9~\%}

{0}{0~\%}

{121}{7.4~\%}

{0}{0~\%}

{490}{8.5~\%}

Ocasional o eventual

{125}{7~\%}

{0}{0~\%}

{180}{7.8~\%}

{0}{0~\%}

{195}{11.9~\%}

{0}{0~\%}

{500}{8.7~\%}

A prueba

{9}{0.5~\%}

{0}{0~\%}

{18}{0.8~\%}

{0}{0~\%}

{10}{0.6~\%}

{0}{0~\%}

{37}{0.6~\%}

Por plazo o tiempodeterminado

{58}{3.2~\%}

{0}{0~\%}

{114}{4.9~\%}

{0}{0~\%}

{111}{6.8~\%}

{0}{0~\%}

{283}{4.9~\%}

Total

{1798}{100~\%}

{0}{100~\%}

{2306}{100~\%}

{0}{100~\%}

{1640}{100~\%}

{0}{100~\%}

{5744}{100~\%}

\hypertarget{tipo-de-contrato}{%
\subsubsection{Tipo de Contrato}\label{tipo-de-contrato}}

\begin{itemize}
\tightlist
\item
  m10. En su trabajo principal, ¿tiene contrato de trabajo escrito?
\end{itemize}

Frecuencias y porcentajes del Tipo de contrato, por ola

Tipo de contrato

Ola

Total

2016

2017

2018

2019

2021

2022

Sí, lo firmé

{1167}{65~\%}

{0}{0~\%}

{1445}{62.7~\%}

{0}{0~\%}

{999}{60.9~\%}

{0}{0~\%}

{3611}{62.9~\%}

Sí, pero no lo hefirmado

{18}{1~\%}

{0}{0~\%}

{25}{1.1~\%}

{0}{0~\%}

{18}{1.1~\%}

{0}{0~\%}

{61}{1.1~\%}

No tengo

{610}{34~\%}

{0}{0~\%}

{833}{36.2~\%}

{0}{0~\%}

{624}{38~\%}

{0}{0~\%}

{2067}{36~\%}

Total

{1795}{100~\%}

{0}{100~\%}

{2303}{100~\%}

{0}{100~\%}

{1641}{100~\%}

{0}{100~\%}

{5739}{100~\%}

\hypertarget{valoraciuxf3n-movimiento-de-demandas-laborales}{%
\subsection{Valoración Movimiento de Demandas Laborales}\label{valoraciuxf3n-movimiento-de-demandas-laborales}}

Considerando la pregunta de ELSOC:

\begin{itemize}
\tightlist
\item
  c20. Pensando en la lista de movimientos sociales que a continuación le mostraré, por favor indique ¿cuál es el que usted más valora?
\end{itemize}

Si consideramos \textbf{solamente} a las personas que responden que el movimiento que más valoran es el de ``Movimiento social de apoyo a demandas laborales'' tenemos:

\begin{table}

\caption{\label{tab:unnamed-chunk-11}Frecuencia de respuestas 'Movimiento social de apoyo a demandas laborales'}
\centering
\begin{tabular}[t]{r|r|r|r}
\hline
ola & n\_mov\_laboral & n\_ola & prop\\
\hline
1 & 381 & 2927 & 0.1301674\\
\hline
2 & 271 & 2473 & 0.1095835\\
\hline
3 & 228 & 3748 & 0.0608324\\
\hline
4 & 311 & 3417 & 0.0910155\\
\hline
5 & 459 & 2740 & 0.1675182\\
\hline
6 & 254 & 2730 & 0.0930403\\
\hline
\end{tabular}
\end{table}

El porcentaje más bajo de respuestas positivas se dio el año 2018 con un 6.08\% y el más alto el año 2021 con un 16.75\%. En promedio corresponde al 10.87\% de la muestra.

Esta variable tiene una serie de ítems asociados al nivel de acuerdo con afirmaciones respecto al movimiento social que más se valora.

\begin{itemize}
\item
  c21. Pensando en {[}movimiento social que más valora{]}, ¿cuán de acuerdo o en desacuerdo está usted con las siguientes afirmaciones?
\item
  c21\_01. Siento un compromiso con este movimiento
\end{itemize}

Frecuencias y porcentajes del Nivel de Compromiso con el Movimiento de Demandas Laborales, por ola

Compromiso con elmovimiento DL

Ola

Total

2016

2017

2018

2019

2021

2022

Totalmente endesacuerdo

{4}{1.1~\%}

{1}{0.4~\%}

{3}{1.3~\%}

{0}{0~\%}

{0}{0~\%}

{4}{1.6~\%}

{12}{0.8~\%}

En desacuerdo

{24}{6.4~\%}

{20}{7.4~\%}

{20}{8.8~\%}

{20}{6.4~\%}

{0}{0~\%}

{22}{8.7~\%}

{106}{7.4~\%}

Ni en desacuerdo nide acuerdo

{41}{10.9~\%}

{33}{12.2~\%}

{30}{13.3~\%}

{31}{10~\%}

{0}{0~\%}

{38}{15~\%}

{173}{12~\%}

De acuerdo

{278}{73.7~\%}

{185}{68.3~\%}

{141}{62.4~\%}

{230}{74~\%}

{0}{0~\%}

{171}{67.6~\%}

{1005}{69.9~\%}

Totalmente deacuerdo

{30}{8~\%}

{32}{11.8~\%}

{32}{14.2~\%}

{30}{9.6~\%}

{0}{0~\%}

{18}{7.1~\%}

{142}{9.9~\%}

Total

{377}{100~\%}

{271}{100~\%}

{226}{100~\%}

{311}{100~\%}

{0}{100~\%}

{253}{100~\%}

{1438}{100~\%}

\begin{itemize}
\tightlist
\item
  c21\_02. Me identifico con este movimiento
\end{itemize}

Frecuencias y porcentajes del Nivel de Identificación con el Movimiento de Demandas Laborales, por ola

Identificación conel movimiento DL

Ola

Total

2016

2017

2018

2019

2021

2022

Totalmente endesacuerdo

{7}{1.8~\%}

{1}{0.4~\%}

{4}{1.8~\%}

{2}{0.6~\%}

{0}{0~\%}

{4}{1.6~\%}

{18}{1.2~\%}

En desacuerdo

{18}{4.7~\%}

{16}{5.9~\%}

{14}{6.2~\%}

{12}{3.9~\%}

{0}{0~\%}

{24}{9.4~\%}

{84}{5.8~\%}

Ni en desacuerdo nide acuerdo

{35}{9.2~\%}

{36}{13.3~\%}

{20}{8.8~\%}

{31}{10~\%}

{0}{0~\%}

{35}{13.8~\%}

{157}{10.9~\%}

De acuerdo

{298}{78.6~\%}

{188}{69.4~\%}

{161}{70.9~\%}

{239}{76.8~\%}

{0}{0~\%}

{176}{69.3~\%}

{1062}{73.6~\%}

Totalmente deacuerdo

{21}{5.5~\%}

{30}{11.1~\%}

{28}{12.3~\%}

{27}{8.7~\%}

{0}{0~\%}

{15}{5.9~\%}

{121}{8.4~\%}

Total

{379}{100~\%}

{271}{100~\%}

{227}{100~\%}

{311}{100~\%}

{0}{100~\%}

{254}{100~\%}

{1442}{100~\%}

\begin{itemize}
\tightlist
\item
  c21\_05. Las acciones y protestas de este movimiento pueden generar un cambio social
\end{itemize}

Frecuencias y porcentajes del Nivel de Acuerdo con acciones del Movimiento de Demandas Laborales generan cambio social, por ola

Acciones movimientoDL pueden generarcambio social

Ola

Total

2016

2017

2018

2019

2021

2022

Totalmente endesacuerdo

{4}{1.1~\%}

{2}{0.7~\%}

{1}{0.4~\%}

{1}{0.3~\%}

{0}{0~\%}

{1}{0.4~\%}

{9}{0.6~\%}

En desacuerdo

{18}{4.7~\%}

{18}{6.6~\%}

{10}{4.4~\%}

{15}{4.9~\%}

{0}{0~\%}

{20}{7.9~\%}

{81}{5.6~\%}

Ni en desacuerdo nide acuerdo

{46}{12.1~\%}

{42}{15.5~\%}

{27}{11.9~\%}

{22}{7.2~\%}

{0}{0~\%}

{31}{12.3~\%}

{168}{11.7~\%}

De acuerdo

{276}{72.8~\%}

{179}{66.1~\%}

{153}{67.4~\%}

{214}{69.7~\%}

{0}{0~\%}

{172}{68.3~\%}

{994}{69.2~\%}

Totalmente deacuerdo

{35}{9.2~\%}

{30}{11.1~\%}

{36}{15.9~\%}

{55}{17.9~\%}

{0}{0~\%}

{28}{11.1~\%}

{184}{12.8~\%}

Total

{379}{100~\%}

{271}{100~\%}

{227}{100~\%}

{307}{100~\%}

{0}{100~\%}

{252}{100~\%}

{1436}{100~\%}

\begin{itemize}
\tightlist
\item
  c21\_09. Las autoridades y el movimiento tienen posiciones opuestas
\end{itemize}

\begin{table}

\caption{\label{tab:unnamed-chunk-15}Frecuencias de 'Grado de acuerdo: Posición opuesta entre el movimiento y las autoridades' por ola}
\centering
\begin{tabular}[t]{l|r|r|r|r|r|r}
\hline
opuesta\_mov & Ola\_1 & Ola\_2 & Ola\_3 & Ola\_4 & Ola\_5 & Ola\_6\\
\hline
Totalmente en desacuerdo & 6 & 3 & 2 & 1 &  & 2\\
\hline
En desacuerdo & 27 & 21 & 14 & 20 &  & 35\\
\hline
Ni en desacuerdo ni de acuerdo & 52 & 39 & 29 & 34 &  & 61\\
\hline
De acuerdo & 250 & 169 & 145 & 211 &  & 144\\
\hline
Totalmente de acuerdo & 37 & 33 & 30 & 37 &  & 6\\
\hline
 & 9 & 6 & 8 & 8 & 459 & 6\\
\hline
\end{tabular}
\end{table}

\begin{itemize}
\tightlist
\item
  c22. Ahora, pensando en los últimos 12 meses, cuán frecuntemente ha participado usted en {[}movimiento social que más valora{]}
\end{itemize}

\begin{table}

\caption{\label{tab:unnamed-chunk-16}Frecuencias de 'Frecuencia: Participación en el movimiento' por ola}
\centering
\begin{tabular}[t]{l|r|r|r|r|r|r}
\hline
participacion\_mov & Ola\_1 & Ola\_2 & Ola\_3 & Ola\_4 & Ola\_5 & Ola\_6\\
\hline
Nunca & 251 & 169 & 148 & 180 & 305 & 164\\
\hline
Casi nunca & 44 & 30 & 22 & 42 & 65 & 38\\
\hline
A veces & 57 & 54 & 38 & 71 & 69 & 37\\
\hline
Frecuentemente & 21 & 15 & 12 & 15 & 19 & 12\\
\hline
Muy Frecuentemente & 6 & 3 & 8 & 3 & 1 & 2\\
\hline
 & 2 &  &  &  &  & 1\\
\hline
\end{tabular}
\end{table}

\hypertarget{participaciuxf3n-en-sindicatos}{%
\subsection{Participación en sindicatos}\label{participaciuxf3n-en-sindicatos}}

\begin{itemize}
\tightlist
\item
  c12\_04. Membresía sindicatos
\end{itemize}

\begin{table}

\caption{\label{tab:unnamed-chunk-17}Frecuencias de cada categoría de Membresía Sindicatos por ola considerando toda la muestra}
\centering
\begin{tabular}[t]{l|r|r|r|r|r|r}
\hline
mem\_sindicato & Ola\_1 & Ola\_2 & Ola\_3 & Ola\_4 & Ola\_5 & Ola\_6\\
\hline
No es miembro & 2613 &  & 3373 &  &  & 2329\\
\hline
Miembro inactivo & 89 &  & 108 &  &  & 48\\
\hline
Miembro activo & 214 &  & 258 &  &  & 193\\
\hline
 & 11 & 2473 & 9 & 3417 & 2740 & 160\\
\hline
\end{tabular}
\end{table}

\begin{table}

\caption{\label{tab:unnamed-chunk-18}Frecuencias de Membresía Sindicato por ola, para quienes valoran movimientos laborales}
\centering
\begin{tabular}[t]{l|r|r|r|r|r|r}
\hline
mem\_sindicato & Ola\_1 & Ola\_2 & Ola\_3 & Ola\_4 & Ola\_5 & Ola\_6\\
\hline
No es miembro & 314 &  & 173 &  &  & 203\\
\hline
Miembro inactivo & 12 &  & 16 &  &  & 7\\
\hline
Miembro activo & 55 &  & 39 &  &  & 25\\
\hline
 &  & 271 &  & 311 & 459 & 19\\
\hline
\end{tabular}
\end{table}

\hypertarget{participaciuxf3n-en-huelgas}{%
\subsection{Participación en huelgas}\label{participaciuxf3n-en-huelgas}}

\begin{itemize}
\tightlist
\item
  c08\_03. Durante los últimos 12 meses, con cuánta frecuencia usted ha participado en una huelga
\end{itemize}

\begin{table}

\caption{\label{tab:unnamed-chunk-19}Frecuencias de cada categoría de Participación en huelga por ola considerando toda la muestra}
\centering
\begin{tabular}[t]{l|r|r|r|r|r|r}
\hline
part\_huelga & Ola\_1 & Ola\_2 & Ola\_3 & Ola\_4 & Ola\_5 & Ola\_6\\
\hline
Nunca & 2547 & 2276 & 3407 & 3069 &  & 2594\\
\hline
Casi nunca & 122 & 78 & 152 & 135 &  & 74\\
\hline
A veces & 188 & 87 & 146 & 134 &  & 49\\
\hline
Frecuentemente & 49 & 25 & 29 & 54 &  & 11\\
\hline
Muy Frecuentemente & 15 & 7 & 11 & 22 &  & 2\\
\hline
 & 6 &  & 3 & 3 & 2740 & \\
\hline
\end{tabular}
\end{table}

\begin{table}

\caption{\label{tab:unnamed-chunk-20}Frecuencias de Participación en huelga por ola, para quienes valoran movimientos laborales}
\centering
\begin{tabular}[t]{l|r|r|r|r|r|r}
\hline
part\_huelga & Ola\_1 & Ola\_2 & Ola\_3 & Ola\_4 & Ola\_5 & Ola\_6\\
\hline
Nunca & 319 & 236 & 196 & 269 &  & 238\\
\hline
Casi nunca & 17 & 12 & 12 & 13 &  & 10\\
\hline
A veces & 33 & 16 & 14 & 24 &  & 4\\
\hline
Frecuentemente & 7 & 6 & 4 & 4 &  & 2\\
\hline
Muy Frecuentemente & 5 & 1 & 2 & 1 &  & \\
\hline
 &  &  &  &  & 459 & \\
\hline
\end{tabular}
\end{table}

\begin{table}

\caption{\label{tab:unnamed-chunk-21}Frecuencias de Participación en huelga por ola, para quienes son miembros de sindicatos (activos o inactivos)}
\centering
\begin{tabular}[t]{l|r|r|r}
\hline
part\_huelga & Ola\_1 & Ola\_3 & Ola\_6\\
\hline
Nunca & 205 & 259 & 198\\
\hline
Casi nunca & 25 & 46 & 23\\
\hline
A veces & 57 & 47 & 11\\
\hline
Frecuentemente & 11 & 8 & 8\\
\hline
Muy Frecuentemente & 5 & 6 & 1\\
\hline
\end{tabular}
\end{table}

\hypertarget{interuxe9s-en-poluxedtica}{%
\subsection{Interés en política}\label{interuxe9s-en-poluxedtica}}

\begin{itemize}
\tightlist
\item
  c13. ¿Qué tan interesado está usted en la política?
\end{itemize}

\begin{table}

\caption{\label{tab:unnamed-chunk-22}Frecuencias de cada categoría de Interés en la política por ola considerando toda la muestra}
\centering
\begin{tabular}[t]{l|r|r|r|r|r|r}
\hline
int\_politica & Ola\_1 & Ola\_2 & Ola\_3 & Ola\_4 & Ola\_5 & Ola\_6\\
\hline
Nada interesado & 1782 & 1269 & 2046 & 1602 & 1108 & 1381\\
\hline
Poco interesado & 509 & 507 & 753 & 674 & 713 & 520\\
\hline
Algo interesado & 365 & 379 & 563 & 601 & 532 & 472\\
\hline
Bastante interesado & 178 & 200 & 264 & 358 & 255 & 258\\
\hline
Muy interesado & 80 & 109 & 110 & 179 & 123 & 99\\
\hline
 & 13 & 9 & 12 & 3 & 9 & \\
\hline
\end{tabular}
\end{table}

\begin{table}

\caption{\label{tab:unnamed-chunk-23}Frecuencias de Interés en la política por ola, para quienes valoran movimientos laborales}
\centering
\begin{tabular}[t]{l|r|r|r|r|r|r}
\hline
int\_politica & Ola\_1 & Ola\_2 & Ola\_3 & Ola\_4 & Ola\_5 & Ola\_6\\
\hline
Nada interesado & 220 & 123 & 121 & 156 & 197 & 108\\
\hline
Poco interesado & 62 & 63 & 44 & 70 & 129 & 59\\
\hline
Algo interesado & 66 & 52 & 38 & 53 & 83 & 52\\
\hline
Bastante interesado & 21 & 20 & 16 & 27 & 34 & 26\\
\hline
Muy interesado & 12 & 12 & 9 & 5 & 16 & 9\\
\hline
 &  & 1 &  &  &  & \\
\hline
\end{tabular}
\end{table}

\begin{table}

\caption{\label{tab:unnamed-chunk-24}Frecuencias de Interés en la política por ola, para quienes son miembros de sindicatos (activos o inactivos)}
\centering
\begin{tabular}[t]{l|r|r|r}
\hline
int\_politica & Ola\_1 & Ola\_3 & Ola\_6\\
\hline
Nada interesado & 158 & 173 & 78\\
\hline
Poco interesado & 49 & 79 & 55\\
\hline
Algo interesado & 55 & 62 & 51\\
\hline
Bastante interesado & 27 & 30 & 40\\
\hline
Muy interesado & 14 & 20 & 17\\
\hline
 &  & 2 & \\
\hline
\end{tabular}
\end{table}

\end{document}
